\documentclass[11pt]{article}

\usepackage[margin=1in]{geometry}
\usepackage{amsmath,amssymb,amsfonts}
\usepackage{booktabs}
\usepackage{siunitx}
\usepackage{graphicx}
\usepackage[round,authoryear]{natbib}
\usepackage[colorlinks=true,linkcolor=blue,citecolor=blue,urlcolor=blue]{hyperref}

\title{AutoRL for Predictive Maintenance on the Edge:\\
REINFORCE with Attention Mechanisms for Milling Tool Replacement}
\author{\textit{Author names omitted for review}}
\date{March 2026}

\begin{document}
\maketitle

\begin{abstract}
Industry~4.0 predictive maintenance (PdM) increasingly runs near machines, where edge constraints on memory, latency, and energy favor compact learning pipelines. This work presents and analyzes an implemented AutoRL pipeline for milling-tool replacement, comparing \textsc{REINFORCE} with \textsc{A2C}, \textsc{DQN}, and \textsc{PPO} under four attention-based feature extractors (Nadaraya--Watson, Temporal, Multi-Head, and Self-Attention). The implementation is grounded in the accompanying codebase (\texttt{rl\_pdm.py}, \texttt{train\_agent.py}): hidden-wear observations, a safety-asymmetric reward function, algorithm-attention sweeps, metadata tracking, and multi-round evaluation with a prognostic $\lambda$ metric. Across 48 trained agents (3 training files $\times$ 4 algorithms $\times$ 4 attentions), evaluated on six unseen datasets over 20 rounds, \textsc{REINFORCE} with temporal or multi-head attention achieves top aggregate performance and statistically outperforms non-\textsc{REINFORCE} baselines on evaluation-score distributions (one-tailed Welch test, $p \ll 0.05$). Results indicate that lightweight policy-gradient learning, when combined with suitable representation bias, can match or exceed heavier optimizers while aligning with edge and sustainability objectives.
\end{abstract}

\textbf{Keywords:} Predictive maintenance; Industry~4.0; Edge computing; AutoRL; \textsc{REINFORCE}; \textsc{PPO}; Attention mechanisms; Sustainable AI.

\section{Context and Motivation}
Predictive maintenance is a core Industry~4.0 capability where machine telemetry is used to schedule interventions before quality loss or failure. Recent reviews emphasize that PdM deployments must balance prediction quality, downtime, and operational constraints \citep{Zonta2020PdMReview,Siraskar2023RLReview}. In industrial environments, edge-centric deployment is often required for low-latency decisions and resilience to network limitations \citep{EdgeIndustrialInternet2021,EdgeComputingIndustrialMachines2019}.

These practical constraints motivate algorithmic efficiency: lower model complexity and fewer optimization passes can reduce training/inference energy and carbon footprint \citep{Strubell2019EnergyPolicy,Patterson2021Carbon,Schwartz2020GreenAI}. This is directly relevant for AutoRL workflows, where many candidate agents are trained and evaluated.

The central question in this manuscript is therefore implementation-oriented: \emph{can a comparatively simple policy-gradient method, \textsc{REINFORCE}, remain competitive against stronger baselines for PdM when combined with attention-based feature extraction?}

\section{AutoRL Pipeline Overview}
The implemented AutoRL workflow automates training, model discovery, and evaluation across algorithms and attention mechanisms. In the provided pipeline:
\begin{itemize}
  \item training is launched via the CLI (\texttt{train\_agent.py}) over algorithm-attention combinations,
  \item each trained model stores metadata (algorithm, attention, learning rate, discount factor, episodes),
  \item evaluation can be single-round or multi-round (20 rounds used here), and
  \item aggregate analytics include heatmaps, grouped comparisons, and statistical testing.
\end{itemize}
This structure follows the broader AutoRL objective of reducing manual tuning burden and improving experimental reproducibility \citep{Faust2019EvolvingRewards,ARLO2023}.

\section{Problem Formulation: MDP vs POMDP}
The maintenance decision process is formulated as an episodic control problem with action space
\begin{equation}
\mathcal{A}=\{0,1\},\quad 0=\text{replace},\; 1=\text{continue}.
\end{equation}
The optimization objective is standard discounted return maximization \citep{SuttonBarto2018}:
\begin{equation}
J(\theta)=\mathbb{E}_{\pi_\theta}\!\left[\sum_{t=0}^{T}\gamma^t r_t\right].
\end{equation}

Although often described using MDP notation, the implemented environment is more precisely a POMDP: true wear is latent to the policy, while observations are sensor-only vectors. In code, \texttt{MT\_Env} excludes \texttt{tool\_wear}, time, and label/action columns from observations, forcing inference from multi-sensor data.

\subsection{Reward and Termination (Code-Faithful)}
With wear threshold $W$, wear ratio $\rho_t=w_t/W$, and constants $(R_1,R_2,R_3,R_4)$ from \texttt{rl\_pdm.py}, the implemented reward is:
\begin{align}
r_t=
\begin{cases}
R_3 + R_4 & a_t=0,\; w_t\le W,\; \rho_t>0.9 \\
R_3 + 0.6R_4 & a_t=0,\; w_t\le W,\; 0.8<\rho_t\le0.9 \\
R_3 + 0.3R_4 & a_t=0,\; w_t\le W,\; 0.7<\rho_t\le0.8 \\
R_3 & a_t=0,\; w_t\le W,\; \rho_t\le0.7 \\
0.5R_2 & a_t=0,\; w_t>W \\
R_2 & a_t=1,\; w_t>W \\
R_1 & a_t=1,\; w_t\le W.
\end{cases}
\end{align}
Episodes terminate on replacement, threshold violation, or end-of-file truncation. This creates a safety-asymmetric objective where violations are much worse than early-but-safe replacement.

\section{Attention Mechanisms for Sensor Representation}
Attention is used here as a feature-extractor layer within RL policy/value networks, rather than as sequence-to-sequence modeling. The general mechanism follows \citet{Vaswani2017Attention}:
\begin{equation}
\mathrm{Att}(Q,K,V)=\mathrm{softmax}\!\left(\frac{QK^\top}{\sqrt{d_k}}\right)V.
\end{equation}

\subsection{Nadaraya--Watson (NW)}
NW attention uses kernel-weighted prototype aggregation \citep{Nadaraya1964,Watson1964}. In implementation, a learnable query projection, learnable keys/values, and positive bandwidth $\beta$ produce
\begin{equation}
\alpha_i=\mathrm{softmax}\left(-\beta\lVert q-k_i\rVert_2^2\right),\qquad
c=\sum_i \alpha_i v_i.
\end{equation}
Its primary advantage is low computational overhead, consistent with edge-oriented PdM findings \citep{Siraskar2024NWAttention}.

\subsection{Temporal (TP)}
Temporal attention in the implemented extractor uses a finite recent history buffer to preserve short-term dynamics across sensor vectors. A learned scorer computes relevance over buffered steps and forms a context vector fused with current features. This is intentionally compact compared with recurrent or full transformer stacks.

\subsection{Multi-Head (MH)}
The Multi-Head extractor projects features into multiple subspaces, applies per-head attention, and recombines outputs. This encourages learning heterogeneous cross-sensor relationships \citep{Vaswani2017Attention} while keeping the overall architecture moderate in size.

\subsection{Self-Attention (SA)}
The Self-Attention extractor provides compact Q/K/V-style interaction over fixed-dimensional observations, with residual and normalization operations for stability. It is more expressive than simple MLP-style extraction but generally heavier than NW/TP.

\section{RL Algorithms Studied}
The study compares \textsc{DQN} \citep{Mnih2015DQN}, \textsc{A2C} \citep{Mnih2016A3C}, \textsc{PPO} \citep{Schulman2017PPO}, and \textsc{REINFORCE} \citep{Williams1992REINFORCE}.

The custom \textsc{REINFORCE} implementation in \texttt{rl\_pdm.py} is built on Stable-Baselines3 on-policy scaffolding and supports an optional value-function baseline (enabled by default) for variance reduction:
\begin{equation}
\nabla_\theta J(\theta)\propto
\sum_t \nabla_\theta\log\pi_\theta(a_t|s_t)\,(G_t-b_t),
\end{equation}
with $b_t=V_\phi(s_t)$ when baseline is used. This detail is important: the implementation is not a strictly baseline-free textbook variant, and this likely contributes to training stability.

\section{Design and Implementation}
\subsection{Environment and Data Interface}
\begin{itemize}
  \item \textbf{Schema-aware features:} \texttt{MT\_Env} detects IEEE/SIT schemas and selects sensor columns accordingly.
  \item \textbf{Hidden wear:} \texttt{tool\_wear} is excluded from observations to enforce inference from sensors.
  \item \textbf{Binary action protocol:} replacement immediately terminates an episode; continuing can trigger catastrophic violation if threshold is crossed.
\end{itemize}

\subsection{Training Orchestration}
\begin{itemize}
  \item \textbf{Grided sweeps:} \texttt{train\_agent.py} supports algorithm, attention, learning-rate, and gamma combinations.
  \item \textbf{Attention mapping:} CLI shortcuts map to implemented extractor types (NW/TP/MH/SA).
  \item \textbf{Metadata persistence:} each agent writes JSON metadata for downstream analysis and reproducibility.
\end{itemize}

\subsection{Evaluation and Scoring}
Evaluation outputs include replacement timing, tool-use percentage, violation statistics, and a prognostic timing metric ($\lambda$) derived from first replacement relative to threshold crossing logic. A weighted evaluation score is computed in \texttt{train\_agent.py} as:
\begin{equation}
\text{EvalScore} = 0.35\cdot \text{ToolUseNorm} + 0.35\cdot (1-\text{ViolNorm}) + 0.30\cdot \Lambda_{\text{norm}},
\end{equation}
where $\Lambda_{\text{norm}}$ penalizes large $|\lambda|$. Multi-round runs (20 rounds) provide distributions suitable for inferential testing.

\section{Experimental Setup}
The reported setup trains 48 agents ($3\times4\times4$) and tests each across six unseen files for 20 rounds. This protocol emphasizes cross-file generalization and robustness to stochasticity. Aggregate comparisons use means, confidence intervals, and one-tailed Welch tests between \textsc{REINFORCE} and pooled non-\textsc{REINFORCE} baselines.

\begin{figure}[t]
\centering
\includegraphics[width=0.95\linewidth]{__Analysis_Report_SIT.png}
\caption{Aggregate analysis report used in the AutoRL comparison pipeline.}
\label{fig:analysisreport}
\end{figure}

\section{Results and Critical Consistency Check}
Results consistently place \textsc{REINFORCE} with TP/MH among top configurations, with high tool-use, low violation tendency, and favorable $\lambda$. The observed significance (very small $p$-values) supports a distributional advantage in this experimental setting.

From a factual review perspective, four points are critical and now aligned to code behavior:
\begin{itemize}
  \item the environment is partially observable (wear hidden),
  \item rewards are threshold-ratio-stepped for replacement timing,
  \item \textsc{REINFORCE} uses an optional baseline (enabled by default), and
  \item evaluation uses a weighted composite rather than raw return alone.
\end{itemize}

These points are central to interpreting why a lightweight optimizer can perform strongly: representation bias and evaluation definition both materially influence ranking.

\section{Discussion}
The main implication is not that \textsc{REINFORCE} is universally superior, but that in this PdM regime (binary action space, strong safety asymmetry, hidden-wear inference), simple policy-gradient updates coupled with appropriate feature extraction can be highly competitive. In AutoRL practice, this is valuable because cheaper algorithms scale better over broad sweep spaces.

This also aligns with Green-AI recommendations to account for computational cost as a first-class concern \citep{Schwartz2020GreenAI,Strubell2019EnergyPolicy}. For edge PdM, the policy class and feature extractor together should be chosen as a compute-performance pair, not independently.

\section{Limitations and Future Work}
Future work should include direct deployment metrics (latency, memory, device-level energy), improved temporal-state reset control for attention buffers at episode boundaries, and robustness tests under sensor drift/domain shift across machines.

\section{Conclusion}
This manuscript integrates implementation-faithful details of the AutoRL codebase with broader scientific context. The combined evidence supports \textsc{REINFORCE}+attention as an effective and resource-aware strategy for edge-oriented milling-tool PdM, while highlighting that evaluation design and representation choices are decisive factors in practical performance.

\section*{Code and Data Availability}
The study is based on the accompanying code and artifacts: \texttt{rl\_pdm.py}, \texttt{train\_agent.py}, and generated evaluation reports.

\begin{thebibliography}{99}

\bibitem[ARLO Team(2023)]{ARLO2023}
M. Mussi, D. Lombarda, A. M. Metelli, F. Trov\`o, and M. Restelli.
ARLO: A framework for automated reinforcement learning.
\emph{Expert Systems with Applications}, 224:119883, 2023.

\bibitem[Carvalho et~al.(2019)]{EdgeComputingIndustrialMachines2019}
A. Carvalho, N. O' Mahony, L. Krpalkova, S. Campbell, J. Walsh, and P. Doody.
Edge computing applied to industrial machines.
\emph{Procedia Manufacturing}, 38:178--185, 2019.

\bibitem[Faust et~al.(2019)]{Faust2019EvolvingRewards}
A. Faust, A. Francis, and D. Mehta.
Evolving rewards to automate reinforcement learning.
\emph{AutoML@ICML}, 2019.

\bibitem[Mnih et~al.(2015)]{Mnih2015DQN}
V. Mnih et~al.
Human-level control through deep reinforcement learning.
\emph{Nature}, 518:529--533, 2015.

\bibitem[Mnih et~al.(2016)]{Mnih2016A3C}
V. Mnih et~al.
Asynchronous methods for deep reinforcement learning.
\emph{Proceedings of ICML}, 2016.

\bibitem[Nadaraya(1964)]{Nadaraya1964}
E. A. Nadaraya.
On estimating regression.
\emph{Theory of Probability and Its Applications}, 9(1):141--142, 1964.

\bibitem[Patterson et~al.(2021)]{Patterson2021Carbon}
D. Patterson, J. Gonzalez, Q. V. Le, et~al.
Carbon emissions and large neural network training.
\emph{arXiv:2104.10350}, 2021.

\bibitem[Schulman et~al.(2017)]{Schulman2017PPO}
J. Schulman, F. Wolski, P. Dhariwal, A. Radford, and O. Klimov.
Proximal policy optimization algorithms.
\emph{arXiv:1707.06347}, 2017.

\bibitem[Schwartz et~al.(2020)]{Schwartz2020GreenAI}
R. Schwartz, J. Dodge, N. A. Smith, and O. Etzioni.
Green AI.
\emph{Communications of the ACM}, 63(12):54--63, 2020.

\bibitem[Siraskar(2025)]{Siraskar2025NaiveREINFORCE}
R. Siraskar.
An empirical study of the na\"{i}ve REINFORCE algorithm for predictive maintenance.
\emph{Discover Applied Sciences}, 2025.

\bibitem[Siraskar et~al.(2023)]{Siraskar2023RLReview}
R. Siraskar, S. Kumar, S. Patil, A. Bongale, and K. Kotecha.
Reinforcement learning for predictive maintenance: a systematic technical review.
\emph{Artificial Intelligence Review}, 2023.

\bibitem[Siraskar et~al.(2024)]{Siraskar2024NWAttention}
R. Siraskar, S. Kumar, S. Patil, A. Bongale, and K. Kotecha.
Application of the Nadaraya--Watson estimator based attention mechanism to predictive maintenance.
\emph{MethodsX}, 12:102754, 2024.

\bibitem[Strubell et~al.(2019)]{Strubell2019EnergyPolicy}
E. Strubell, A. Ganesh, and A. McCallum.
Energy and policy considerations for deep learning in NLP.
\emph{Proceedings of ACL}, 2019.

\bibitem[Sutton and Barto(2018)]{SuttonBarto2018}
R. S. Sutton and A. G. Barto.
\emph{Reinforcement Learning: An Introduction}.
MIT Press, 2nd edition, 2018.

\bibitem[Vaswani et~al.(2017)]{Vaswani2017Attention}
A. Vaswani et~al.
Attention is all you need.
\emph{Advances in Neural Information Processing Systems}, 2017.

\bibitem[Watson(1964)]{Watson1964}
G. S. Watson.
Smooth regression analysis.
\emph{Sankhy\={a}: The Indian Journal of Statistics, Series A}, 1964.

\bibitem[Williams(1992)]{Williams1992REINFORCE}
R. J. Williams.
Simple statistical gradient-following algorithms for connectionist reinforcement learning.
\emph{Machine Learning}, 8:229--256, 1992.

\bibitem[Zhang et~al.(2021)]{EdgeIndustrialInternet2021}
T. Zhang, Y. Li, and C. L. P. Chen.
Edge computing and its role in Industrial Internet: methodologies, applications, and future directions.
\emph{Information Sciences}, 557:34--65, 2021.

\bibitem[Zonta et~al.(2020)]{Zonta2020PdMReview}
T. Zonta et~al.
Predictive maintenance in Industry 4.0: a systematic literature review.
\emph{Computers \& Industrial Engineering}, 150:106889, 2020.

\end{thebibliography}

\end{document}